\centerline{\bf \Huge Векторная арифметика}

\begin{flushright}
        \scriptsize{Если вас съели, у вас есть два выхода}
\end{flushright}

\textbf{Определение 1.} \textit{Направленным отрезком} называется упорядоченная пара точек на плоскости. Направленный отрезок с началом в точке \(A\) и концом в точке \(B\) обозначается \( \overline{AB} \) или \( \overrightarrow{AB} \). Начало и конец направленного отрезка могут совпадать.

\textbf{Определение 2.} Направленные отрезки \(\overrightarrow{AB}\) и \(\overrightarrow{CD}\) называются \textit{эквивалентными}, если \(ABDC\) -- параллелограмм (возможно, вырожденный). Эквивалентность направленных отрезков является отношением эквивалентности.

\textbf{Определение 3.} Классы эквивалентности направленных отрезков называются \textit{векторами}. Словосочетание «вектор \(\overrightarrow{AB}\)» означает класс эквивалентности, содержащий направленный отрезок \(\overrightarrow{AB}\).

\textbf{Определение 4.} Векторы называются \textit{коллинеарными}, если направленные отрезки, соответствующие этим векторам, отложенные от одной точки, лежат на одной прямой. Если они при этом лежат на одном луче, исходящем из этой точки, то они называются \textit{сонаправленными}. Если на разных --- то \textit{противоположно направленными}.

\problem{1} 
а) Упростите выражение \( \overrightarrow{AC} + \overrightarrow{DE} + \overrightarrow{CB} + \overrightarrow{EA} \).

\noindent
б) Упростите выражение \( \overrightarrow{CE} - \overrightarrow{CA} + \overrightarrow{EB} - \overrightarrow{DB} \).

\problem{2}
а) Пусть \(AA_1\) -- медиана треугольника \(ABC\). Докажите, что \( \overrightarrow{AA_1} = \frac{1}{2} (\overrightarrow{AB} + \overrightarrow{AC}) \).

\noindent
б) На стороне \(BC\) треугольника \(ABC\) выбрана точка \(A_1\) такая, что \( \frac{\overline{BA_1}}{\overline{BC}} = t \). Докажите, что \( \overrightarrow{AA_1} = (1 - t) \overrightarrow{AB} + t \overrightarrow{AC} \).

\problem{3}
Пусть на плоскости нарисован четырехугольник $ABCD$ и отмечены середины сторон $AB$ и $CD$ точки $M$ и $N$ соответственно. Докажите, что $\overrightarrow{MN} = \dfrac{1}{2}(\overrightarrow{AD} + \overrightarrow{BC})$

\problem{4}
Точки $M, K, N$ и $L$ $-$ середины сторон $AB, BC, CD$ и $DE$ пятиугольника $ABCDE$ (не обязательно выпуклого), $P$ и $Q \ -$ середины отрезков $MN$ и $KL$. Докажите, что отрезок $PQ$ в четыре раза меньше стороны $AE$ и параллелен ей.

\problem{5}
Пусть $ABCD$ и $AB_1C_1D_1$ $-$ два параллелограмма с общей вершиной. Докажите, что один из векторов $\overrightarrow{BB_1}, \overrightarrow{CC_1}$ и $\overrightarrow{DD_1}$ коллинеарен сумме двух других.




